% Page 072 — page imprimée 68
% La famille Roussel du Villaret au Château d’Ayres (suite)

\textbf{Ginette et Nicole}

\medskip

Ginette Roussel s’était mariée en 1934 ou 1935 avec Jean Teissier du Cros dont la famille,
originaire de Valleraugue, était connue.

\medskip

\textit{(Une des sœurs de Jean, Germaine Dieterlen-Teissier du Cros -ethnologue- devint la collaboratrice de Marcel
Griaule, spécialiste des Dogons, avec qui elle écrivit « Renard pâle », puis l’assistante et compagne de Jean
Rouch ethnologue et cinéaste, auteur de « Moi un Noir »...)}

\medskip

Il avait été en poste en Martinique puis en Chine enfin en Indochine où il emmena sa jeune
femme. Elle n’apprécia sans doute guère la vie « coloniale » ni l’atmosphère de Saïgon. Le
château d’Ayres lui manquait et ils revinrent s’y installer peu de temps avant la guerre. Jean
paraissait de santé délicate, le teint jaune ; il souffrait d’une amibiase chronique dont il ne put
jamais se débarrasser. De pension de familles, le château devint un hôtel réputé sous la
direction ferme de Madame Roussel.

\medskip

Au château était arrivée en 1928, Nicole Roussel qui passait pour être la sœur de Ginette. Elle
était née en 1927, et on disait qu’elle était la fille de Melle Harrison, sœur de Mme. Roussel
qui l’aurait adoptée Elle allait à l’école de filles à Meyrueis et faisait tourner la tête de
plusieurs garçons de l’école contiguë car elle était jolie fille et surtout chantait très bien sans
timidité aucune. Bref, j’en fus amoureux comme quelques autres et plus tard en été 1944 ou
1945 nous nous retrouvâmes dans le groupe de jeunes de la paroisse de Meyrueis pour
quelques balades et réunions diverses.

\medskip

En fait Nicole était, disait la rumeur publique, la fille de Frédéric H*, tailleur à Paris, (ou d’un
de ses amis) qu’il avait eue d’une de ses conquêtes parisiennes venue en séjour au château
d’Ayres. Quelques années après, elle épousa un médecin devint Mme. A.. (Mme A mère, était
une Monod) et passa toute sa vie à Châlons sur Marne où elle mourut il y a quelques années.
Elle avait perdu un fils atteint d’un cancer et elle ne s’en était jamais consolée. Son second fils
serait ingénieur en Allemagne

\medskip

.

C’est Ginette qui hérita du château et qui, avec son mari, mais aussi grâce à son maître d’hôtel
– chef du personnel - continua et améliora, la réputation du château qui accueillit quelques
personnalités politiques célèbres telles Adenauer ou de Gaulle ainsi que des artistes connus…

\medskip

Ce maître d’hôtel avait une réputation détestable mais, doué, il avait su se rendre
indispensable. Plusieurs de ceux qui l’ont connu le qualifient encore de « sale bête » (sic).
Protestant, eh oui ! il portait à la ceinture une grosse croix huguenote ancienne. Elle voisinait
avec l’insigne des SAC : les Services d’Action Civique de Charles Pasqua, qui s’illustrèrent
par de nombreuses exactions et sans doute par des crimes. Il ne cachait pas, bien au contraire,
cet engagement politique d’extrême droite !!! Dans le village les jeunes et les enfants
l’avaient surnommé « Casserole ».

\medskip

* Frédéric H. voir note « in fine »
